\documentclass[11pt]{article}
\usepackage{ochamath}
\subjectcolor{2F5DA8}
\setsheet{線形代数}{行列と線形変換　演習}{https://university-math-crj.pages.dev/linear-algebra/matrices/}
\begin{document}
\makesheethead{解答}

\problem[サイズの判定]
\begin{ans}
$AB$ は内側のサイズ2が一致するので定義でき、$3\times4$ 行列となる。$BA$ は内側が4と3で一致しないため定義できない。
\end{ans}
\point{サイズは計算前に判定できる。}

\problem[積と非可換性]
\begin{ans}
行と列を掛けて足すと \[AB=\begin{pmatrix}4&2\\1&1\end{pmatrix},\qquad BA=\begin{pmatrix}2&4\\1&3\end{pmatrix}.\] また $A(2,-1)^{\mathsf T}=(2-2,-1)^{\mathsf T}=(0,-1)^{\mathsf T}$。例えば積の第1成分は $1\times2+2\times1=4$。
\end{ans}
\point{順番を変えた積も成分から確認する。}
\newpage

\problem[回転]
\begin{ans}
$R(2,1)^{\mathsf T}=(-1,2)^{\mathsf T}$ で反時計回りの90度回転。$R^2=\begin{pmatrix}-1&0\\0&-1\end{pmatrix}=-I$ なので180度回転である。長さは回転前後で $\sqrt5$ のまま。
\end{ans}
\point{回転の合成は角度の和に対応する。}

\problem[センサの出力]
\begin{ans}
\[\boldsymbol y=\begin{pmatrix}1&1\\1&-1\end{pmatrix}\boldsymbol x,\qquad \boldsymbol y=\begin{pmatrix}5+2\\5-2\end{pmatrix}=\begin{pmatrix}7\\3\end{pmatrix}.\] 第1行は和、第2行は差を作る。
\end{ans}
\point{各行が出力1つ分の式になる。}
\end{document}
